\documentclass[11pt]{article}
\usepackage[margin=1in]{geometry}
\usepackage{amsmath,amssymb,amsthm}
\usepackage{mathtools}
\usepackage{enumitem}
\usepackage{hyperref}
\usepackage{cleveref}
\usepackage{booktabs}
\usepackage{tabularx}
\usepackage{microtype}
\usepackage{cite}
\usepackage{xurl}
\usepackage{path}
\hypersetup{colorlinks=true,linkcolor=blue,citecolor=blue,urlcolor=blue}
\newtheorem{definition}{Definition}
\newtheorem{lemma}{Lemma}
\newtheorem{remark}{Remark}
\title{Public Request Notice Selections for Source-Only Releases\\\large Choosing Canonical Notice Families from Within-Release Status and Cross-Release Carryforward Tokens}
\author{Anonymous}
\date{Unified archive note v0.07 (2026-03-21; archive v473)}
\begin{document}
\maketitle

\begin{abstract}
The archive can now say what a source-only paper promises on request, when that promise is complete, whether the paper-level public token survives a successor release, what exact outward-facing sentence readers saw, and what canonical sentence family sits behind that sentence. One smaller publication question still recurs in terse papers and release notes: given the imported status or carryforward basis tokens, which maintained canonical notice family should be selected? This note gives that choice rule one home. It defines compact \emph{public request notice selections}: token-keyed selector objects that map an imported within-release status envelope or successor-release carryforward envelope to one canonical public-request-notice normal form and one emitted notice instance, without re-owning either the narrower status semantics or the sentence family itself.
\end{abstract}

\paragraph{Non-redundancy note.}
Synthesis~59 owns the within-release paper-level public status token, Synthesis~60 owns the cross-release carryforward verdict, Synthesis~61 owns the emitted outward-facing notice sentence and pointer order, Synthesis~62 owns the canonical sentence families and slot bindings behind those notices, Synthesis~64 owns the notice-keyed follow-up response menu used after that family choice is fixed, Synthesis~73 owns the request-side terminal stop later terse papers should usually cite, Synthesis~74--75 own the paired terminal / cutover discipline, and Synthesis~17 owns the concrete worked instantiation.\cite{series:publicrequeststatusenvelopes,series:publicrequestcarryforward,series:publicrequestnotices,series:publicrequestnoticenormalforms,series:publicrequestresponsemenus,series:publicrequestresponsepacketrefreshresponsepacketclosureverdicts,series:pairedterminalcitationdiscipline,series:pairedterminalcutoverrule,series:workedexample}
This note owns only the selector saying which canonical source-only notice family applies to a given imported token combination, and later terse papers should cite it directly only when that exact token-to-family choice is itself live.

\paragraph{Scope.}
This is not a new public request contract, not a new completion criterion, not a new paper-level token, not a new carryforward verdict, not a replacement for the emitted notice, and not a replacement for the canonical sentence family.
It is a choice-rule note for selecting among already-owned source-only public-request notice families.

\paragraph{Exports.}
This note exports (i) public request notice selections, (ii) a selector-totality rule, (iii) a basis-token choice rule, and (iv) an emitted-agreement rule separating the choice of canonical family from the family itself and from the concrete emitted notice. Later terse papers should cite Synthesis~63 only when the live issue is which canonical source-only notice family was selected for one explicit imported token bundle; stop at Synthesis~59--60 when the live issue is the imported within-release status token or successor carryforward verdict, stop at Synthesis~61 when the live issue is the exact emitted outward-facing sentence or pointer order, stop at Synthesis~62 when the live issue is the reusable sentence family or slot template, widen to Synthesis~64 when the live issue is which smallest-sufficient maintained follow-up object should be handed over once that family choice is fixed, widen through Synthesis~65--72 when the live issue has moved into the exact predecessor packet cut or the successor reuse / refresh / wrapper / handoff chain, widen to Synthesis~73 when the live issue is whether that request-side branch already closes, widen to Synthesis~74--75 when the remaining issue is paired-terminal reuse or fail-closed cutover, and stop at Synthesis~17 when the concrete maintained worked route is the live issue instead.\cite{series:publicrequeststatusenvelopes,series:publicrequestcarryforward,series:publicrequestnotices,series:publicrequestnoticenormalforms,series:publicrequestresponsemenus,series:publicrequestresponsepackets,series:publicrequestresponsepacketcarryforward,series:publicrequestresponsepacketdeltaledgers,series:publicrequestresponsepacketrefreshnotices,series:publicrequestresponsepacketrefreshnormalforms,series:publicrequestresponsepacketrefreshselections,series:publicrequestresponsepacketrefreshresponsemenus,series:publicrequestresponsepacketrefreshresponsepackets,series:publicrequestresponsepacketrefreshresponsepacketclosureverdicts,series:pairedterminalcitationdiscipline,series:pairedterminalcutoverrule,series:workedexample}

\section{Why notice selections deserve their own note}
Once a source-only branch has both emitted notices and canonical notice families, later terse papers often start to blur two different questions.
One question is which sentence family is canonical.
Another is which of those maintained families is actually licensed by the basis-token combination present in the imported status or carryforward owner.
Without a separate selector owner, neighboring notes either restate token-to-family choice inline or quietly let the normal-form note absorb a policy decision that belongs one layer later.
A selector keeps that choice auditable without turning the canonical family note into a hidden decision table.

\begin{definition}[Public request notice selection]
Fix one emitted public request notice $N$ and one maintained public request notice normal form $F$ whose basis is either a within-release public request status envelope or a successor-release public request carryforward envelope.
A \emph{public request notice selection} is a tuple
\[
S=(\beta,\iota,\phi,\nu,\sigma)
\]
consisting of one imported basis object $\beta$, one selector-input token bundle $\iota$, one selected normal-form key $\phi$, one selected emitted-notice key $\nu$, and one rendered sentence $\sigma$ that agrees with both the selected normal form and the emitted notice.
The selector owns only the choice of canonical family for that token bundle.
It does not re-own the basis object, the normal-form template, or the emitted notice instance.
\end{definition}

\begin{definition}[Selector-totality rule]
A public request notice selection satisfies the \emph{selector-totality rule} if every emitted source-only public-request notice shipped for the maintained release cut appears in exactly one selector entry whose basis-token bundle is explicit.
So later terse papers do not need to guess whether a notice family choice was left implicit.
\end{definition}

\begin{definition}[Basis-token choice rule]
A public request notice selection satisfies the \emph{basis-token choice rule} if the imported status or carryforward tokens in $\iota$ match the imported basis object and if the selected normal form allows exactly that token combination.
For a within-release selection, the public-request-status token must agree with the imported status envelope.
For a successor-carryforward selection, both the successor public-request-status token and the carryforward token must agree with the imported carryforward envelope.
So the selector may not choose a canonical notice family that the basis object does not license.
\end{definition}

\begin{definition}[Emitted-agreement rule]
A public request notice selection satisfies the \emph{emitted-agreement rule} if the rendered sentence $\sigma$ is exactly the rendered sentence of the selected normal form and exactly the sentence shipped by the selected emitted notice, and if the selector changes neither that notice's reader-pointer order nor its narrower owner artifacts.
So the selector chooses among already-owned notice families; it does not rewrite the shipped sentence.
\end{definition}

\begin{lemma}[Selectors isolate token-to-family choice without replacing family or notice ownership]
If a public request notice selection satisfies the selector-totality, basis-token choice, and emitted-agreement rules, then a later terse paper may cite the selector when it needs to justify why a particular canonical source-only notice family was chosen for a given imported token combination, while still citing the normal form for the family itself and the emitted notice for the exact visible sentence and pointer order.
\end{lemma}

\begin{proof}
The selector-totality rule makes the choice explicit for every shipped notice rather than leaving some choices implicit in prose.
The basis-token choice rule prevents the selector from licensing a normal form unsupported by the imported status or carryforward owner.
The emitted-agreement rule prevents the selector from silently rewriting the sentence or taking over the concrete wrapper instance.
Therefore the selector contributes only the token-to-family choice, leaving family ownership with the normal form and emitted-sentence ownership with the notice.\qedhere
\end{proof}

\begin{remark}[Minimal citation rule]
Cite the public request notice selection when the claim is about \emph{which canonical source-only notice family was chosen for a particular token combination}.
Cite the public request notice normal form when the claim is about \emph{the reusable sentence family or slot template itself}.
Cite the public request notice when the claim is about \emph{which exact outward-facing sentence and pointer order were actually emitted}.
Cite the status or carryforward envelope when the claim is about \emph{which basis token or cross-release verdict is true}.
\end{remark}

\begin{remark}[Selection stop discipline]
When later terse papers need only the stable request-side stop or the stop-versus-widen judgment for the visible source-only branch, they should usually cite Synthesis~73--75 rather than expand this selector layer locally.
Widen through this note only when the exact token-to-family choice is itself disputed. In that case, the first sufficient audit stop is the selector entry itself, and widening should proceed only to the selected normal form, then to the selected emitted notice, and only then to the imported status or carryforward basis object named by that selector.
So the selector owns which family was chosen for one explicit token bundle, while sentence-family ownership, emitted wrapper ownership, and narrower basis semantics remain with their already-maintained neighbors.
\end{remark}

\paragraph{A forced public notice-selector matrix.}
\begin{table}[t]
\centering
\small
\begin{tabularx}{0.97\linewidth}{@{}p{0.31\linewidth}X@{}}
\toprule
\textbf{Forced row} & \textbf{Forced value}\\
\midrule
Within-release selector row & \texttt{public\_request\_notice\_selection\_key = worked\_example\_receipt\_interlock\_within\_release\_public\_request\_notice\_selection} with \texttt{notice\_kind = within\_release}; the only worked selector-input bundle is \texttt{public\_request\_status\_token = complete} with no carryforward token \\
Successor selector row & \texttt{public\_request\_notice\_selection\_key = worked\_example\_receipt\_interlock\_successor\_public\_request\_notice\_selection} with \texttt{notice\_kind = successor\_carryforward}; the only worked selector-input bundle is \texttt{public\_request\_status\_token = complete} plus \texttt{carryforward\_token = preserved\_complete} \\
Chosen-family floor & The within-release row forces \texttt{worked\_example\_receipt\_interlock\_within\_release\_complete\_public\_request\_notice\_normal\_form}; the successor row forces \texttt{worked\_example\_receipt\_interlock\_successor\_preserved\_complete\_public\_request\_notice\_normal\_form} because those are the only worked families whose allowed-token sets match the imported bundles \\
Chosen-wrapper floor & The selected emitted wrappers remain \texttt{worked\_example\_receipt\_interlock\_public\_request\_status\_notice} and \texttt{worked\_example\_receipt\_interlock\_public\_request\_carryforward\_notice}; the selector records which wrapper/family pair is live and does not invent a new sentence \\
Decision-label floor & \texttt{selection\_decision = select\_within\_release\_complete\_notice\_normal\_form} within release and \texttt{selection\_decision = select\_successor\_preserved\_complete\_notice\_normal\_form} across the successor pair \\
Rendered-agreement floor & Each selector row may be quoted only while the chosen family and chosen wrapper still agree on the same rendered sentence; disagreement reopens the family or emitted-wrapper owner rather than leaving the selector unchanged \\
Staleness trigger & This matrix goes stale immediately if the selector-input bundle changes, if the chosen family's allowed-token set no longer matches that bundle, or if the emitted wrapper stops inheriting the same rendered sentence; then the selector must choose differently or fail closed \\
Escalation pointer & Widen to Synthesis~59--60 for basis-token / carryforward semantics, Synthesis~61 for the emitted wrapper, Synthesis~62 for the reusable family itself, and Synthesis~64--75 once the live issue has moved into follow-up handoff, packet maintenance, or request-side terminal / cutover ownership \\
\bottomrule
\end{tabularx}
\caption{A forced public notice-selector matrix. This is the smallest public answer later terse papers can quote directly when the live issue is which canonical source-only notice family was forced for one explicit imported token bundle rather than the narrower basis-token owner, the reusable family definition, or the wider request-side routing tail.}
\label{tab:forced-public-notice-selector-matrix}
\end{table}

This matrix is intentionally non-decorative: it pins the actual worked within-release and successor-facing selector keys, the exact imported selector-input tokens, the forced family/wrapper choices, the explicit decision labels, and the point at which the selector citation has gone stale without silently re-owning the basis tokens, the reusable sentence family, or the emitted wrapper instance.\cite{series:workedexample}

\paragraph{One tiny selector-forcing drill.}
In the worked successor selector, the imported token bundle is exactly \texttt{public\_request\_status\_token = complete} together with \texttt{carryforward\_token = preserved\_complete}. Under that bundle, the only honest worked decision is \texttt{select\_successor\_preserved\_complete\_notice\_normal\_form}, because that is the maintained family whose allowed-token set matches both imported fields and whose rendered sentence still agrees with the emitted wrapper. If the carryforward envelope instead published \texttt{reopened} or \texttt{advanced}, or if the emitted wrapper ceased to realize the same rendered sentence, this matrix could no longer be quoted unchanged: the selector would have to choose a different maintained family, or fail closed until one exists.\cite{series:workedexample}

\section{Worked example pointer}
The maintained worked bundle now ships \nolinkurl{example_public_request_notice_selections.json}.\cite{series:workedexample}
It contains exactly two selector entries keyed \nolinkurl{worked_example_receipt_interlock_within_release_public_request_notice_selection} and \nolinkurl{worked_example_receipt_interlock_successor_public_request_notice_selection}.
The within-release selector imports the explicit token bundle \nolinkurl{public_request_status_token=complete}, chooses normal form \nolinkurl{worked_example_receipt_interlock_within_release_complete_public_request_notice_normal_form}, chooses emitted notice \nolinkurl{worked_example_receipt_interlock_public_request_status_notice}, records the decision label \nolinkurl{select_within_release_complete_notice_normal_form}, and agrees on the exact rendered sentence ``For this source-only release, the worked note's paper-level public request promise is complete and supporting artifacts remain available on request.''
The successor-facing selector imports the explicit token bundle \nolinkurl{public_request_status_token=complete} plus \nolinkurl{carryforward_token=preserved_complete}, chooses normal form \nolinkurl{worked_example_receipt_interlock_successor_preserved_complete_public_request_notice_normal_form}, chooses emitted notice \nolinkurl{worked_example_receipt_interlock_public_request_carryforward_notice}, records the decision label \nolinkurl{select_successor_preserved_complete_notice_normal_form}, and agrees on the exact rendered sentence ``Across the canned successor release, that complete paper-level public request status carries forward unchanged.''
That gives later terse papers one citeable owner for the token-to-family choice without confusing that choice with either the canonical family or the concrete emitted notice. In the maintained worked bundle, those two exact selector keys remain route-visible beside the imported notice and normal-form maps so later terse papers can recover the whole visible source-only wrapper ladder from one maintained route object without restating the token bundle in prose. Later terse papers should usually combine that route-visible selector layer with Synthesis~73--75 rather than narrate the whole visible source-only branch locally.\cite{series:publicrequestresponsepacketrefreshresponsepacketclosureverdicts,series:pairedterminalcitationdiscipline,series:pairedterminalcutoverrule}
Once that choice is fixed, Synthesis~64 owns the next paper-facing question: which smallest sufficient maintained object should be handed over for a concrete follow-up about that visible notice.

\begin{thebibliography}{99}\small

\bibitem{series:publicrequeststatusenvelopes}
Anonymous.
\newblock Public Request Status Envelopes for Source-Only Releases: Paper-Level Tokens Compressing Family Completion without Re-owning the Promise.
\newblock Series synthesis note (Synthesis~59), 2026. (This archive.)

\bibitem{series:publicrequestcarryforward}
Anonymous.
\newblock Public Request Carryforward for Source-Only Successor Releases: When Paper-Level Request Status Persists, Advances, or Reopens across Release Evolution.
\newblock Series synthesis note (Synthesis~60), 2026. (This archive.)

\bibitem{series:publicrequestnotices}
Anonymous.
\newblock Public Request Notices for Source-Only Releases: Outward-Facing Sentences for Within-Release Status and Cross-Release Carryforward.
\newblock Series synthesis note (Synthesis~61), 2026. (This archive.)

\bibitem{series:publicrequestnoticenormalforms}
Anonymous.
\newblock Public Request Notice Normal Forms for Source-Only Releases: Canonical Sentence Families for Within-Release Status and Cross-Release Carryforward Notices.
\newblock Series synthesis note (Synthesis~62), 2026. (This archive.)

\bibitem{series:publicrequestresponsemenus}
Anonymous.
\newblock Public Request Response Menus for Source-Only Releases: Smallest-Sufficient Follow-Up Handoffs for Basis Tokens, Visible Notices, Canonical Families, Family Choice, Scope, and Completion.
\newblock Series synthesis note (Synthesis~64), 2026. (This archive.)

\bibitem{series:publicrequestresponsepackets}
Anonymous.
\newblock Public Request Response Packets for Source-Only Releases: Exact Maintained Handoff Slices for Concrete Follow-Up Challenges.
\newblock Series synthesis note (Synthesis~65), 2026. (This archive.)

\bibitem{series:publicrequestresponsepacketcarryforward}
Anonymous.
\newblock Public Request Response Packet Carryforward for Source-Only Successor Releases: When Exact Handoff Slices Survive, Refresh, or Reopen.
\newblock Series synthesis note (Synthesis~66), 2026. (This archive.)

\bibitem{series:publicrequestresponsepacketdeltaledgers}
Anonymous.
\newblock Public Request Response Packet Delta Ledgers for Source-Only Successor Releases: Exact Packet-Local Refresh Maps across Successor Evolution.
\newblock Series synthesis note (Synthesis~67), 2026. (This archive.)

\bibitem{series:publicrequestresponsepacketrefreshnotices}
Anonymous.
\newblock Public Request Response Packet Refresh Notices for Source-Only Successor Releases: Visible Reissue Sentences for Refreshed Handoff Slices.
\newblock Series synthesis note (Synthesis~68), 2026. (This archive.)

\bibitem{series:publicrequestresponsepacketrefreshnormalforms}
Anonymous.
\newblock Public Request Response Packet Refresh Notice Normal Forms for Source-Only Successor Releases: Canonical Visible Families for Refreshed Handoff Slices.
\newblock Series synthesis note (Synthesis~69), 2026. (This archive.)

\bibitem{series:publicrequestresponsepacketrefreshselections}
Anonymous.
\newblock Public Request Response Packet Refresh Notice Selections for Source-Only Successor Releases: Choosing Canonical Visible Families from Packet-Local Refresh Tokens.
\newblock Series synthesis note (Synthesis~70), 2026. (This archive.)

\bibitem{series:publicrequestresponsepacketrefreshresponsemenus}
Anonymous.
\newblock Public Request Response Packet Refresh Response Menus for Source-Only Successor Releases: Smallest-Sufficient Successor-Facing Handoffs after Visible Refresh.
\newblock Series synthesis note (Synthesis~71), 2026. (This archive.)

\bibitem{series:publicrequestresponsepacketrefreshresponsepackets}
Anonymous.
\newblock Public Request Response Packet Refresh Response Packets for Source-Only Successor Releases: Exact Successor-Facing Handoff Slices after Visible Refresh.
\newblock Series synthesis note (Synthesis~72), 2026. (This archive.)

\bibitem{series:publicrequestresponsepacketrefreshresponsepacketclosureverdicts}
Anonymous.
\newblock Public Request Response Packet Refresh Response Packet Closure Verdicts for Source-Only Successor Releases: Capstone Certificates for Exact Successor-Facing Handoff Slices.
\newblock Series synthesis note (Synthesis~73), 2026. (This archive.)

\bibitem{series:pairedterminalcitationdiscipline}
Anonymous.
\newblock Paired Terminal Citation Discipline for Stable Review Spines: One Answer-Side Stop and One Request-Side Capstone for Terse Papers.
\newblock Series synthesis note (Synthesis~74), 2026. (This archive.)

\bibitem{series:pairedterminalcutoverrule}
Anonymous.
\newblock Paired Terminal Cutover Rules for Stable Review Spines: When Later Terse Papers Must Stop, Widen, or Reopen.
\newblock Series synthesis note (Synthesis~75), 2026. (This archive.)

\bibitem{series:workedexample}
Anonymous.
\newblock Worked Example: Receipt Interlock for an Anonymous-DHT Reader-Privacy Claim.
\newblock Series synthesis note (Synthesis~17), 2026. (This archive.)

\end{thebibliography}
\end{document}
